\clearpage
\section{Considerații teoretice}

Dezvoltarea produselor soft fără o proiectare preventivă a început să devină
imposibilă oradă cu creșterea dimensiunilor produselor soft dezvoltate. Astfel
au fost create mai multe limbaje de modelare care aveau ca scop eliminarea
acestei probleme. Unele limbaje de modelare obiect orientate au început să
apare în perioada între mijlocul anilor 1970 şi sfîrşitul anilor 1980, cînd
diferiţi cercetători şi programatori propuneau diverse metode de analiza şi
proiectare obiect orientată (APOO). În perioada între anii 1989-1994 numărul
total al limbajelor de modelare cele mai cunoscute a crescut de la 10 pîna la
mai mult decît 50. Mulţi utilizatori  au avut dificultăţi  serioase la alegerea
limbajului  de APOO, fiindcă  nici unul din limbajele propuse nu satisfăcea
toate cerinţele către modelarea (construirea modelelor) sistemelor compuse. \cite{Melnic} În
cele din urmă mai multe organizații au creat UML care continuă să se dezvolte
și să fie utilizat până în prezent.

\subsection{Limbajul UML}
UML, prescurtarea de la Unified Modeling Language, este un limbaj de modelare
standardizat format dintr-un set integrat de diagrame, dezvoltat pentru a ajuta
dezvoltatorii de sisteme și software pentru specificarea, vizualizarea,
construirea și documentarea artefactelor sistemelor software, precum și pentru
modelarea afacerilor și a altor sisteme non-software. UML reprezintă o colecție
de cele mai bune practici de inginerie care s-au dovedit a fi de succes în
modelarea sistemelor mari și complexe. UML este o parte foarte importantă a
dezvoltării software-ului orientat pe obiecte și a procesului de dezvoltare
software. UML folosește mai ales notații grafice pentru a exprima proiectarea
proiectelor software. Utilizarea UML ajută echipele de proiect să comunice, să
exploreze proiectele potențiale și să valideze designul arhitectural al
software-ului. În acest articol, vă vom oferi idei detaliate despre ceea ce
este UML, Istoricul UML și o descriere a fiecărui tip de diagramă UML, împreună
cu exemple UML. \cite{UML}

UML definește treisprezece tipuri de diagrame, împărțite în trei categorii:
șase tipuri de diagrame reprezintă structura statică a aplicației; trei
reprezintă tipuri generale de comportament; și patru reprezintă aspecte
diferite ale interacțiunilor:
\begin{itemize}
\item \textbf{Diagramele de structură} includ diagrama de clasă, diagrama obiectului, diagrama componentelor, diagrama structurii compozite, diagrama pachetului și diagrama de implementare; 
\item \textbf{Diagramele de comportament} includ diagrama cazului de Utilizare (utilizată de unele metodologii în timpul Adunării cerințelor); diagrama de activitate și diagrama mașinii de stare; 
\item \textbf{Diagramele de interacțiune}, toate derivate din diagrama de comportament mai generală, includ diagrama de secvență, diagrama de comunicare, diagrama de sincronizare și diagrama de prezentare a interacțiunii.
\end{itemize}

Limbajul UML este bazat pe un anumit număr de noţiuni principale care pot fi
studiate şi aplicate de către majoritatea programiştilor şi elaboratorilor
cunoscuţi cu metodele de analiza şi proiectarea obiect orientate. Totodată
noţiunile de bază pot fi combinate şi extinse în asa fel că specialiştii
modelării orientate pe obiecte pot elabora de sinestătător modele ale
sistemelor complexe în diferite domenii de aplicare.


\vfill
\subsubsection{Diagrama cazului de utilizare}

Modelarea vizuală în UML poate fi reprezentată ca un oarecare proces de lansare
pe niveluri de la cel mai general şi abstract model conceptual al sistemului
iniţial către model logic şi mai apoi fizic, ce corespunde unui sistem de
program. Pentru atingerea acestui scop de la început se crează un model în
formă de diagarama cazurilor de utilizare (use case diagram) care descrie
destinaţia functională a sistemului sau cu alte cuvinte descrie ceea ce
sistemul va executa în procesul său de funcţionare. Diagrama cazurilor de
utilizare reprezintă un model iniţial conceptual al unui sistem în procesul de
proiectare şi exploatare.

Proiectarea a unei diagrame a cazurilor de utilizare urmăreşte scopurile următoare:
\begin{itemize}
\item determinarea limitelor comune şi a contextului domeniului de modelare la
    etapele iniţiale de proiectare a unui sistem;

\item formularea cerinţelor comune către comportare funcţională a sistemului
    proiectat;

\item elaborarea modelului iniţial conceptual al unui sistem pentru detalierea
    de mai tîrziu în forma modelelor logice şi fizice;

\item pregătirea documentărei iniţiale pentru interacţiunea elaboratorilor unui
    sitem cu clienţii şi utilizatorii.

\end{itemize}

În figura \ref{fig:ex_usecase} este prezentat exemplul unei diagrame de caz de utilizare.
\begin{figure}[H]
    \begin{center}
        \includegraphics[width=15cm]{ex/usecase}
        \caption{Exemplu diagramă caz de utilizare \cite{usecase}}
        \label{fig:ex_usecase}
    \end{center}
\end{figure}


\vfill
\subsubsection{Diagrama secvență}

În limbajul UML colaborarea între entităţi se cercetează în aspectul informativ
al comunicaţiilor lor, adică obiectele care interacţionează între ele fac un
oarecare schimb de informaţii. Pentru modelarea colaborării între obiecte în
limbajul UML se utilizează diagramele de interacţiune. Vorbind despre aceste
diagrame se iau în consideraţie două aspecte.

Mai  întîi,  colaborarea  între  obiecte  poate  fi  cercetată  în timp  şi
atunci  pentru  reprezentarea particularităţilor temporale şi modului de
acceptare a mesajelor se utilizează diagrama de secvenţă.

În al doilea rînd pot fi cercetate particularităţile structurale ale
colaborării între obiecte. Pentru reprezentarea particularităţilor de
transmitere şi acceptare a mesajelor între obiecte se utilizează diagrama de
colaborare. 

În figura \ref{fig:ex_secvence} este prezentat exemplul unei diagrame de secvență.
\begin{figure}[H]
    \begin{center}
        \includegraphics[width=15cm]{ex/secvence}
        \caption{Exemplu diagramă de secvență \cite{secvence}}
        \label{fig:ex_secvence}
    \end{center}
\end{figure}

\vfill
\subsubsection{Diagrama de colaborare}

Particularitatea principală a diagramei de colaborare constă în posibilitatea de a reprezenta grafic nu
numai consecutivitatea colaborării dar şi toate relaţii structurale între obiecte. În primul rînd în
diagrama de colaborare sub formă de dreptunghiuri se reprezintă obiectele care conţin numele
obiectului, clasei, valorile atributului. Mai departe, ca şi în diagrama de clase se indică asocierile
între obiecte sub forma de linii de conectare. Totodată pot fi indicate numele asocierilor şi al
rolurilor obiectelor pentru asocierea dată. Suplimentar pot fi reprezentate legăturile dinamice –
fluxurile de mesaje. Ele se reprezintă ca linii între obiecte cu săgeţi care indică derecţia, numele
mesajului şi numărul de ordine în consecutivitatea de iniţializare a mesajelor.

Spre deosebire de diagrama de secvenţă în diagrama de colaborare sunt reprezentate relaţiile între
obiecte care sunt importante pentru colaborare. Din altă parte în această diagrama nu se indică
timpul în calitate de măsură. Consecutivitatea de acţiuni şi flux paralel pot fi determinate cu ajutorul
numerelor de ordine, deci este posibilă specificarea legăturilor între obiecte în timp real.

Pentru atingerea unui scop sau pentru realizarea unui serviciu comportamentul unui sistem poate fi
descris la nivelul obiectelor care fac schimb de mesaje. Din punct de vedere a unui analist sau a
unui constructor este importantă reflectarea legăturilor structurale ale obiectelor aparte. Diagrama
de  colaborare  reflectă  un  fel  de  reprezentare  statică  a  unui  sistem  ca  totalitate  de  obiecte
dependente.

În figura \ref{fig:ex_collaboration} este prezentat exemplul unei diagrame de colaborare.
\begin{figure}[H]
    \begin{center}
        \includegraphics[width=15cm]{ex/colaborare}
        \caption{Exemplu diagramă de colaborare \cite{collaboration}}
        \label{fig:ex_collaboration}
    \end{center}
\end{figure}

\vfill
\subsubsection{Diagramă de clasă}

Locul central în APOO îl ocupă elaborarea modelului logic al unui sistem în
forma diagramei de clase. Notaţia claselor în limbajul UML este simplă şi clară
pentru toţi. O notaţie asemănătoare se utilizează şi pentru obiecte – care sunt
exemplare ale clasei, unica diferenţă este în aceea că la numele clasei se
adaugă numele obiectului şi toată înregistrarea se subliniază.

Notaţia limbajului UML oferă multe posibilităţi pentru reflectarea informaţiei
suplimentare (operaţii abstracte  sau  clase,  stereotipuri,  metode  comune
şi  particulare,  interfeţe  detaliate,  clase parametrizate). Totodată este
posibilă utilizarea reprezentării grafice pentru asocieri şi proprietăţile lor
specifice, astfel cum sunt relaţiile de agregare, cînd ca părţi componente ale
clasei pot fi alte clase.

Diagrama de clase (class diagram) se utilizează pentru reprezentarea structurii
statice a unui model de sistem în terminologia claselor programării OO.
Diagrama de clase poate reflecta diferite legături între entităţile domeniului
de obiecte (obiecte şi subsisteme) şi descrie structura lor internă şi tipurile
de relaţii. În această diagramă nu este menţionată informaţia despre aspectele
temporare ale funcţionării sistemului. Din acest punct de vedere diagrama de
clase este dezvoltarea ulterioară a

Diagrama de clase reprezintă un graf cu noduri – elemente de tip
«clasificatori» care sunt legate prin diferite tipuri de relaţii de structură.
Trebuie de menţionat că diagrama de clase poate conţine interfeţe, pachete,
relaţii şi chiar exemplare, aşa ca obiecte şi legături. Prin diagrama de clase
se subînţelege modelul structural static al sistemului proiectat, de accea
diagrama de clase deseori se socoate o reprezentare grafică a legăturilor
structurale ale modelului logic al sistemului care sunt independente şi
invariante în timp.

În figura \ref{fig:ex_class} este prezentat exemplul unei diagrame de clasă.
\begin{figure}[H]
    \begin{center}
        \includegraphics[width=15cm]{ex/clasa}
        \caption{Exemplu diagramă de clasă \cite{class}}
        \label{fig:ex_class}
    \end{center}
\end{figure}

\vfill
\subsubsection{Diagrama de stare}

Pentru modelarea comporatamentului la nivelul logic în limbajul UML pot fi
utilizate la rînd cîteva diagrame canonice: de stări, de activitate, de
secvenţă şi de cooperare, fiecare din care pune accentul  la  aspectul
specificat  de funcţionare  a  sistemului.  Spre deosebire  de alte  diagrame,
diagrama de stări descrie procesul de modificare a stărilor numai pentru o
clasă, pentru un exemplar a unei clase, adică de a modela toate modificările
posibile în starea unui  propriu obiect. În urma căruia modificarea stării
obiectului poate fi provocată de influenţa externă a altor obiecte sau din
exterior. Anume pentru descrierea reacţiei obiectelor la aşa fel de influenţe
externe sunt utilizate diagramele de stări.

Această diagramă este folosită pentru descrierea consecutivităţilor de stări
posibile şi trecerilor, care în ansamblu caracterizează comportamentul
elementelor modelului în timpul ciclului de viaţă. Diagrama de stări reprezintă
comportamentul dinamic a entităţilor în baza specificaţiei reacţiei lor la
perceperea căror-va evenimente concrete.

Deşi diagramele de stări sunt mai rar utilizate pentru descrierea
comporatamenrului de care-va exemplare  a  claselor  (obiectelor),  dar ele  de
asemenea  pot fi  utilizate  şi pentru  specificarea funcţionalităţilor altor
componente a modelului, aşa cum cazurile de utilizare, actorii, subsisteme,
operaţii şi metode.

Diagrama  de stări în realitate este un graf de înfăţişare specială, care
reprezintă un automat. Definiţia de automat în contextul UML are o semantică
prea specifică, bazată pe teoria automatelor. Vîrfurile acestui graf sunt
stările şi alte elemente a automatului (pseudostări), care sunt reprezentate cu
ajutorul simbolilor grafice speciale. Razele grafului sunt pentru marcarea
trecerilor de la o stare la altă. Diagramele de stări pot fi depuse una în
alta, formînd diagramele depuse de reprezentarea mai detaliată a căror-va
elemente a modelului. Pentru înţelegerea semanticii diagramei de stări concrete
este necesar de a imagina nu numai deosebirile în comportamentul entităţii
modelate, dar şi de a şti noţiuni generale din teoria automatelor.

În figura \ref{fig:ex_state} este prezentat exemplul unei diagrame de stare.
\begin{figure}[H]
    \begin{center}
        \includegraphics[width=15cm]{ex/stare}
        \caption{Exemplu diagramă de stare \cite{state}}
        \label{fig:ex_state}
    \end{center}
\end{figure}


\vfill
\subsubsection{Diagrama de activități}

Pentru modelarea procesului de executare a operaţiilor în limbajul UML se
utilizează aşa numitele diagrame de activităţi. Notaţia grafica acceptată
pentru aceste diagrame are mult comun cu notaţia diagramei de stări ce se
evidenţiază prin notarea stărilor şi tranziţiilor. Deosebirea constă în
semantica stărilor care sunt utilizate pentru prezentarea acţiunilor dar nu
activităţilor, şi în aceea că tranziţiile evenimentelor nu sunt etichetate.
Fiecare stare în diagrama de activităţi corespunde executării unei operaţiuni
elimentare, dar trecere în altă stare se execută numai la terminarea operaţiei
în starea precedentă. Grafic diagrama de activităţi se reprezintă în forma unui
graf de activitate cu nodurile – stări activitate şi muchile – tranziţii de la
o stare la altă.

În aşa fel, diagramele de activităţi pot fi considerate cazuri particulare ale
diagramelor de stări. Anume aceste diagrame permit realizarea administrării
procedurale şi sincrone care depinde de terminarea activităţii interne în
limbajul UML. Sensul principal al utilizării acestor diagrame constă în
vizualizarea particularităţilor operaţiilor unor clase pentru reprezentarea
algoritmilor executării lor. Totodată fiecare stare realizează operaţiile unei
clase anumite şi permite utilizarea diagramei de activităţi pentru descrierea
reacţiilor la evenimente interne acestui sistem.

În contextul limbajului UML activitatea (activity) reprezintă o totalitate de
calcule executate de către automat. Totodată calculele elementare pot duce la
un anumit rezultat sau careva acţiune (action). În diagrama de activităţi se
reflectă logica sau consecutivitatea tranziţiilor de la o acţiune la alta,
totodată se evidenţiază rezultatul activităţii. Rezultatul, la rîndul său poate
duce la schimbarea stării sistemului dat sau la returnarea unei valori.

În figura \ref{fig:ex_activity} este prezentat exemplul unei diagrame de activitate.
\begin{figure}[H]
    \begin{center}
        \includegraphics[width=15cm]{ex/activitate}
        \caption{Exemplu diagramă de activitate \cite{activity}}
        \label{fig:ex_activity}
    \end{center}
\end{figure}

\subsection{Programarea orientata pe obiect}

Pentru a putea implementa cât mai fidel sistemul modelat
este necesar să utilizăm un limbaj care suportă POO.
Acest lucru nu este o restricție, POO fiind considerat un pas important în evoluția limbajelor de programare.
POO apărut din necesitatea exprimării problemei într-un mod mai natural
ființei umane. Astfel unitățile care alcătuiesc un program se apropie mai mult
de modul nostru de a gîndi decît modul de lucru al calculatorului. Pînă la
apariția programării orientate pe obiect, programele erau implementate în
limbaje de programare procedural(C, Pascal) sau în limbaje care nici măcar nu
oferea o madalitate de gruparea instrucțiunilor în unități logice, precum
funcții sau proceduri cum este cazul limbajului de asamblare Assembly.


Principalele principii ale programării orientate pe obiect sunt:
\begin{itemize}
    \item încapsularea;
    \item abstracția;
    \item moștenirea;
    \item polimorfizmul.
\end{itemize}


\subsubsection{ Încapsularea }
Să zicem că avem un program. Are câteva obiecte logic diferite care comunică
între ele — conform regulilor definite în program. Încapsularea se realizează
atunci când fiecare obiect își păstrează starea privată , în interiorul unei
clase. Alte obiecte nu au acces direct la această stare. În schimb, ei pot
apela doar o listă de funcții publice - numite metode. Deci, obiectul își
gestionează propria stare prin metode - și nicio altă clasă nu îl poate atinge
decât dacă este permis în mod explicit. Dacă doriți să comunicați cu obiectul,
ar trebui să utilizați metodele furnizate. Dar (în mod implicit), nu puteți
schimba starea. Să presupunem că construim un mic joc Sims. Sunt oameni și
există o pisică. Ei comunică între ei. Dorim să aplicăm încapsularea, așa că
încapsulăm toată logica „pisica” într-o Cat clasă.  Puteți hrăni pisica. Dar nu
poți schimba direct cât de foame este pisica. Aici „starea” pisicii este
variabilele mood private hungryși energy. Are și o metodă privată meow(). Îl
poate suna oricând dorește, celelalte clase nu pot spune pisicii când să
miaună. Ceea ce pot face ei este definit în metodele sleep() publice play()și
feed(). Fiecare dintre ele modifică cumva starea internă și poate invoca
meow(). Astfel, se face legătura între statul privat și metodele publice.
Aceasta este încapsularea.

\subsubsection{ Abstracția }
Abstracția poate fi gândită ca o extensie naturală a încapsulării. În
proiectarea orientată pe obiecte, programele sunt adesea extrem de mari. Și
obiectele separate comunică foarte mult între ele. Așa că menținerea unei baze
de cod mare ca aceasta de ani de zile - cu modificări pe parcurs - este
dificilă. Abstracția este un concept care urmărește să ușureze această
problemă. Aplicarea abstracției înseamnă că fiecare obiect ar trebui să expună
doar un mecanism de nivel înalt pentru utilizarea acestuia. Acest mecanism ar
trebui să ascundă detaliile interne de implementare. Ar trebui să dezvăluie
doar operațiuni relevante pentru celelalte obiecte. Gândiți-vă - un aparat de
cafea. Face o mulțime de lucruri și scoate zgomote ciudate sub capotă. Dar tot
ce trebuie să faci este să pui cafea și să apeși un buton. De preferință, acest
mecanism ar trebui să fie ușor de utilizat și să se schimbe rareori în timp.
Gândiți-vă la el ca la un mic set de metode publice pe care orice altă clasă le
poate apela fără să „știe” cum funcționează. Un alt exemplu de abstractizare
din viața reală? Gândiți-vă la modul în care vă folosiți telefonul: Telefoanele
mobile sunt complexe. Dar folosirea lor este simplă. Interacționați cu
telefonul folosind doar câteva butoane. Ce se întâmplă sub capotă? Nu trebuie
să știți — detaliile de implementare sunt ascunse. Trebuie doar să cunoașteți
un set scurt de acțiuni. Modificările de implementare — de exemplu, o
actualizare de software — afectează rareori abstracția pe care o utilizați.

\subsubsection{ Moștenirea }
OK, am văzut cum încapsularea și abstractizarea ne pot ajuta să dezvoltăm și să
menținem o bază de cod mare. Dar știți care este o altă problemă comună în
proiectarea OOP? Obiectele sunt adesea foarte asemănătoare. Ei împărtășesc o
logică comună. Dar nu sunt în totalitate la fel. Uf… Deci, cum reutilizam
logica comună și extragem logica unică într-o clasă separată? O modalitate de a
realiza acest lucru este moștenirea. Înseamnă că creați o clasă (copil)
derivând dintr-o altă clasă (părinte). Astfel, formăm o ierarhie. Clasa copil
reutiliza toate câmpurile și metodele clasei părinte (partea comună) și își
poate implementa propriile sale (partea unică). Un profesor privat este un tip
de profesor. Și orice profesor este un tip de Persoană. Dacă programul nostru
trebuie să gestioneze profesori publici și privați, dar și alte tipuri de
oameni, cum ar fi studenții, putem implementa această ierarhie de clasă. În
acest fel, fiecare clasă adaugă doar ceea ce este necesar pentru ea în timp ce
reutiliza logica comună cu clasele părinte.

\subsubsection{ Polimorfizmul }
Suntem până la cel mai complex cuvânt! Polimorfismul înseamnă „multe forme” în
greacă. Deci știm deja puterea moștenirii și o folosim cu bucurie. Dar vine
această problemă. Să presupunem că avem o clasă părinte și câteva clase copil
care moștenesc de la ea. Uneori dorim să folosim o colecție – de exemplu o
listă – care conține o combinație a tuturor acestor clase. Sau avem o metodă
implementată pentru clasa părinte – dar am dori să o folosim și pentru copii.
Acest lucru poate fi rezolvat folosind polimorfismul. Mai simplu spus,
polimorfismul oferă o modalitate de a utiliza o clasă exact ca părintele ei,
astfel încât să nu existe confuzie cu tipurile de amestecare. Dar fiecare clasă
de copii își păstrează propriile metode așa cum sunt. Acest lucru se întâmplă
de obicei prin definirea unei interfețe (părinte) care să fie reutilizată.
Prezintă o grămadă de metode comune. Apoi, fiecare clasă copil implementează
propria sa versiune a acestor metode. De fiecare dată când o colecție (cum ar
fi o listă) sau o metodă așteaptă o instanță a părintelui (unde sunt subliniate
metode comune), limbajul are grijă de evaluarea implementării corecte a metodei
comune - indiferent de care copil este trecut. Aruncă o privire la o schiță a
implementării figurilor geometrice.  Triunghiul, Cercul și Dreptunghiul pot fi
folosite acum în aceeași colecție Având aceste trei cifre care moștenesc
părintele, Figure Interfacevă permite să creați o listă de triangles, circles,
și rectangles. Și tratați-i ca pe același tip de obiect. Apoi, dacă această
listă încearcă să calculeze suprafața pentru un element, se găsește și se
execută metoda corectă. Dacă elementul este un triunghi, CalculateSurface() se
numește triunghiul. Dacă este un cerc, atunci CalculateSurface() se numește
cerc. Si asa mai departe. Dacă aveți o funcție care operează cu o figură
folosind parametrul acesteia, nu trebuie să o definiți de trei ori - o dată
pentru un triunghi, un cerc și un dreptunghi. Îl poți defini o dată și accepta
a Figure ca argument. Indiferent dacă treceți un triunghi, un cerc sau un
dreptunghi - atâta timp cât implementează CalculateParamter(), tipul lor nu
contează. Sper că acest lucru a ajutat. Puteți folosi direct aceleași
explicații la interviurile de angajare. Dacă găsiți ceva încă greu de înțeles -
nu ezitați să întrebați în comentariile de mai jos.
